\rowcolors{0}{}{}
\arrayrulecolor{black}
\renewcommand{\arraystretch}{1.3}
\small
\begin{longtable}{|C{1cm}|L{5.6cm}|L{8cm}|C{0.8cm}|}
\hline
\textbf{ID US} & \textbf{User Story} & \textbf{Description technique principale} & \textbf{Effort} \\
\hline
\endfirsthead

\hline

\endhead

\hline
\endfoot

\hline
\endlastfoot

% Epic 8: Gestion des demandes clients
\multicolumn{4}{|l|}{\cellcolor{blue!20}\textbf{Epic 8 - Gestion des demandes clients}} \\
\hline

US38 & En tant que \textbf{client}, je souhaite soumettre une demande de maintenance afin de signaler un besoin. &
\begin{itemize}[nosep, left=0pt, label={$\bullet$}]
\item Développer l'endpoint \footnotesize{"POST /api/demandes"}\small.
\item Valider l'éligibilité de la demande (contrat en cours) et calculer automatiquement l'échéance de prise en charge (SLA).
\item Assurer le téléversement de pièces jointes justificatives.
\item Réaliser l'interface de dépôt de dossier côté Next.js.
\end{itemize} & 3 \\
\hline

US39 & En tant que \textbf{chef de projet}, je souhaite consulter les demandes en attente afin de les traiter. &
\begin{itemize}[nosep, left=0pt, label={$\bullet$}]
\item Structurer un flux d'extraction filtré selon le périmètre opérationnel du chef de projet connecté.
\item Mettre en place des tris multicritères (antériorité, urgence) et concevoir une interface adaptative (liste, calendrier ou Kanban).
\end{itemize} & 2 \\
\hline

US40 & En tant que \textbf{chef de projet}, je souhaite valider une demande afin de la convertir en ticket. &
\begin{itemize}[nosep, left=0pt, label={$\bullet$}]
\item Développer l'endpoint \footnotesize{"PATCH /api/demandes/:id/valider"}\small.
\item Exécuter la routine de conversion automatique de la demande validée en ticket d'intervention et notifier le client.
\end{itemize} & 2 \\
\hline

US41 & En tant que \textbf{chef de projet}, je souhaite demander des informations complémentaires afin de clarifier une demande. &
\begin{itemize}[nosep, left=0pt, label={$\bullet$}]
\item Développer l'endpoint \footnotesize{"PATCH /api/demandes/:id/demander-infos"}\small.
\item Basculer le statut de la demande en attente, consigner l'action à l'historique et alerter le client.
\end{itemize} & 2 \\
\hline

US42 & En tant que \textbf{client}, je souhaite compléter ma demande afin de fournir les informations manquantes. &
\begin{itemize}[nosep, left=0pt, label={$\bullet$}]
\item Développer l'endpoint \footnotesize{"PATCH /api/demandes/:id/fournir-infos"}\small{} recueillant les réponses du client.
\item Réactualiser l'état d'avancement de la demande et notifier le chef de projet référent.
\end{itemize} & 2 \\
\hline

US43 & En tant que \textbf{chef de projet}, je souhaite requalifier une demande afin de corriger sa catégorie ou sa priorité. &
\begin{itemize}[nosep, left=0pt, label={$\bullet$}]
\item Développer l'endpoint \footnotesize{"PATCH /api/demandes/:id/requalifier"}\small{} avec obligation de renseigner un motif.
\item Inscrire les modifications typologiques au journal d'audit pour garantir la cohérence des indicateurs.
\end{itemize} & 1 \\
\hline

US44 & En tant qu'\textbf{utilisateur}, je souhaite commenter une demande afin de communiquer sur le contexte. &
\begin{itemize}[nosep, left=0pt, label={$\bullet$}]
\item Développer l'endpoint \footnotesize{"POST /api/demandes/:id/commentaires"}\small{} pour centraliser le fil de discussion.
\item Assurer l'historisation des échanges et l'envoi d'alertes aux acteurs engagés sur la demande.
\end{itemize} & 2 \\
\hline

% Epic 9: Gestion des tickets
\multicolumn{4}{|l|}{\cellcolor{blue!20}\textbf{Epic 9 - Gestion des tickets}} \\
\hline

US46 & En tant que \textbf{Chef de Projet ou Administrateur}, je souhaite visualiser et gérer l'ensemble des tickets dans un tableau Kanban afin de modifier leurs statuts et d'organiser efficacement le flux de travail. &
\begin{itemize}[nosep, left=0pt, label={$\bullet$}]
\item Développer l'endpoint \footnotesize{"GET /api/tickets/kanban"}\small{} ségréguant les tickets par colonnes d'état.
\item Modéliser l'interface Kanban interactive en Next.js avec gestion des transitions fluides et affichage des alertes de SLA.
\end{itemize} & 3 \\
\hline

US47 & En tant que \textbf{chef de projet}, je souhaite assigner un ticket à un développeur afin de répartir la charge. &
\begin{itemize}[nosep, left=0pt, label={$\bullet$}]
\item Développer les endpoints d'affectation unitaire (\footnotesize{"POST /api/tickets/:id/assign"}\small) et de récupération des collaborateurs disponibles.
\item Garantir l'unicité de l'affectation principale et expédier une notification au développeur retenu.
\end{itemize} & 2 \\
\hline

US48 & En tant que \textbf{chef de projet}, je souhaite assigner un ticket à plusieurs développeurs afin de collaborer sur des tâches complexes. &
\begin{itemize}[nosep, left=0pt, label={$\bullet$}]
\item Étendre l'endpoint d'assignation pour supporter une équipe de co-résolution sur un même ticket.
\item Présenter un aperçu de la charge de travail planifiée sur 7 jours pour chaque profil disponible.
\item Journaliser le collectif d'intervention dans l'historique d'audit du ticket.
\end{itemize} & 2 \\
\hline

US49 & En tant que \textbf{développeur}, je souhaite me désassigner d'un ticket afin de libérer ma charge. &
\begin{itemize}[nosep, left=0pt, label={$\bullet$}]
\item Développer l'endpoint \footnotesize{"POST /api/tickets/:id/unassign"}\small{} conditionné par la saisie obligatoire d'une justification.
\item Notifier l'encadrement technique et actualiser instantanément le journal d'activité du ticket.
\end{itemize} & 1 \\
\hline

US50 & En tant que \textbf{développeur}, je souhaite démarrer un chronomètre sur un ticket afin de tracker mon temps. &
\begin{itemize}[nosep, left=0pt, label={$\bullet$}]
\item Développer l'endpoint \footnotesize{"POST /api/tickets/chrono/lancer/:assignmentId"}\small{} pour instancier une session d'activité active.
\item Bloquer toute tentative d'activation si une autre session de suivi du temps est déjà en cours pour l'utilisateur.
\item Réaliser l'affichage du widget de chronométrage en temps réel sur le bandeau Next.js.
\end{itemize} & 2 \\
\hline

US51 & En tant que \textbf{développeur}, je souhaite mettre en pause mon chronomètre afin de gérer les interruptions. &
\begin{itemize}[nosep, left=0pt, label={$\bullet$}]
\item Développer l'endpoint \footnotesize{"POST /api/tickets/chrono/pause/:assignmentId"}\small{} lié à un référentiel de motifs de suspension.
\item Figer l'état de la session courante tout en sauvegardant le cumul des minutes consommées.
\end{itemize} & 1 \\
\hline

US52 & En tant que \textbf{développeur}, je souhaite arrêter le chronomètre afin de finaliser ma session. &
\begin{itemize}[nosep, left=0pt, label={$\bullet$}]
\item Développer l'endpoint \footnotesize{"POST /api/tickets/chrono/arreter/:assignmentId"}\small{} consolidant la session de travail.
\item Clôturer la mesure, stocker la durée finale sommée et inscrire la ligne de temps dans l'historique de production.
\end{itemize} & 1 \\
\hline

US53 & En tant que \textbf{développeur}, je souhaite ajouter manuellement une session de travail afin de corriger des oublis. &
\begin{itemize}[nosep, left=0pt, label={$\bullet$}]
\item Développer l'endpoint \footnotesize{"POST /api/tickets/chrono/manuel"}\small{} permettant la saisie a posteriori d'un volume horaire.
\item Valider la cohérence des dates saisies par rapport au cycle de vie et aux affectations du ticket.
\end{itemize} & 2 \\
\hline

US54 & En tant que \textbf{développeur}, je souhaite changer le statut d'un ticket afin de refléter l'avancement. &
\begin{itemize}[nosep, left=0pt, label={$\bullet$}]
\item Développer l'endpoint \footnotesize{"PATCH /api/tickets/:id/statut"}\small{} contrôlant les règles de transition d'états.
\item Auditer le changement d'état et pousser des notifications automatisées aux parties prenantes.
\end{itemize} & 2 \\
\hline

US55 & En tant que \textbf{client}, je souhaite valider la résolution d'un ticket afin de confirmer que le problème est résolu. &
\begin{itemize}[nosep, left=0pt, label={$\bullet$}]
\item Développer l'endpoint \footnotesize{"POST /api/tickets/:id/validate-resolution"}\small{} actant la fin de l'intervention.
\item Recueillir l'avis de clôture du client, verrouiller le ticket et envoyer une notification de succès à l'équipe technique.
\end{itemize} & 1 \\
\hline

US56 & En tant que \textbf{chef de projet}, je souhaite rouvrir un ticket fermé afin de traiter un problème récurrent. &
\begin{itemize}[nosep, left=0pt, label={$\bullet$}]
\item Développer l'endpoint \footnotesize{"POST /api/tickets/:id/signaler-probleme-non-resolu"}\small{} suite à une réclamation.
\item Réinitialiser le cycle de vie du ticket vers un état d'analyse et ré-alerter la maîtrise d'œuvre.
\end{itemize} & 1 \\
\hline

US57 & En tant que \textbf{chef de projet}, je souhaite consulter le temps passé par ticket afin d'analyser la productivité. &
\begin{itemize}[nosep, left=0pt, label={$\bullet$}]
\item Développer un service d'agrégation calculant le ratio entre le temps estimé initial et le cumul des sessions réelles.
\item Modéliser les graphiques de consommation horaire directement dans la vue détaillée du ticket.
\end{itemize} & 2 \\
\hline

US58 & En tant que \textbf{développeur}, je souhaite signaler un blocage afin d'alerter sur un obstacle. &
\begin{itemize}[nosep, left=0pt, label={$\bullet$}]
\item Mettre en place la procédure de levée d'incident bloquant sur une affectation avec description obligatoire de la cause.
\item Taguer visuellement le ticket comme "bloqué" et émettre une alerte prioritaire vers l'encadrement.
\end{itemize} & 2 \\
\hline

US59 & En tant que \textbf{membre d'équipe}, je souhaite résoudre un blocage afin de permettre au développeur de reprendre son travail. &
\begin{itemize}[nosep, left=0pt, label={$\bullet$}]
\item Développer la cinématique d'arbitrage (soumission, validation ou rejet d'une solution de déblocage).
\item Restaurer le statut de production nominal du ticket et archiver la résolution pour la base de connaissances.
\end{itemize} & 1 \\
\hline

US60 & En tant que \textbf{développeur}, je souhaite demander une extension de délai SLA afin de disposer de plus de temps. &
\begin{itemize}[nosep, left=0pt, label={$\bullet$}]
\item Développer l'endpoint \footnotesize{"POST /api/tickets/:id/request-extension"}\small{} motivé par des contraintes techniques avérées.
\item Transmettre la requête d'avenant au responsable de projet et acter l'attente dans la timeline.
\end{itemize} & 1 \\
\hline

US61 & En tant que \textbf{chef de projet}, je souhaite approuver ou rejeter les demandes d'extension SLA. &
\begin{itemize}[nosep, left=0pt, label={$\bullet$}]
\item Développer les interfaces de décision d'extension (validation ou refus avec recalcul automatique de la date butoir du SLA).
\item Notifier le développeur et consigner la modification de la métrique contractuelle.
\end{itemize} & 1 \\
\hline

US62 & En tant qu'\textbf{utilisateur}, je souhaite commenter un ticket avec mentions afin de communiquer efficacement. &
\begin{itemize}[nosep, left=0pt, label={$\bullet$}]
\item Développer l'endpoint \footnotesize{"POST /api/tickets/:id/commentaires"}\small{} prenant en charge l'identification ciblée d'utilisateurs.
\item Segmenter l'affichage pour isoler les notes internes (visibilité équipe) des commentaires partagés (visibles par le client).
\end{itemize} & 2 \\
\hline

US63 & En tant qu'\textbf{utilisateur}, je souhaite ajouter des pièces jointes aux tickets afin de fournir du contexte. &
\begin{itemize}[nosep, left=0pt, label={$\bullet$}]
\item Développer l'endpoint \footnotesize{"POST /api/tickets/:id/commentaires/with-attachments"}\small.
\item Assurer le contrôle de conformité des fichiers (poids, extensions autorisées) et leur association croisée au fil de discussion.
\end{itemize} & 2 \\
\hline

US64 & En tant que \textbf{développeur}, je souhaite être informé lorsqu'un autre utilisateur modifie le même ticket. &
\begin{itemize}[nosep, left=0pt, label={$\bullet$}]
\item Mettre en œuvre un service de synchronisation d'état en temps réel pour signaler la présence ou l'édition simultanée.
\item Rafraîchir dynamiquement les données à l'écran pour prévenir les collisions et écrasements de saisie.
\end{itemize} & 3 \\
\hline

US65 & En tant qu'\textbf{utilisateur}, je souhaite consulter l'historique complet d'un ticket afin de comprendre son évolution. &
\begin{itemize}[nosep, left=0pt, label={$\bullet$}]
\item Développer l'endpoint \footnotesize{"GET /api/tickets/:id/history"}\small{} compilant l'ensemble des mutations de l'entité.
\item Restituer une frise chronologique (timeline) unifiée mêlant changements de statuts, assignations successives et jalons de vie.
\end{itemize} & 2 \\
\hline

% Epic 16: Assistance intelligente
\multicolumn{4}{|l|}{\cellcolor{blue!20}\textbf{Epic 16 - Assistance intelligente (IA)}} \\
\hline

US103 & En tant que \textbf{chef de projet}, je souhaite que l'IA analyse automatiquement les blocages afin d'identifier rapidement les problèmes critiques. &
\begin{itemize}[nosep, left=0pt, label={$\bullet$}]
\item S'appuyer sur le service d'analyse des blocages exposé par l'API IA pour générer un diagnostic à partir du ticket signalé.
\item Restituer une synthèse du blocage avec un niveau de confiance et l'afficher dans l'interface de suivi.
\end{itemize} & 3 \\
\hline

US104 & En tant que \textbf{chef de projet}, je souhaite que l'IA propose des actions correctives afin de débloquer rapidement un développeur. &
\begin{itemize}[nosep, left=0pt, label={$\bullet$}]
\item Exploiter les endpoints d'analyse de blocage de l'API IA pour proposer des pistes de résolution à partir du contexte du ticket.
\item Permettre au chef de projet de consulter les suggestions et de suivre l'escalade ou la validation de la solution proposée.
\end{itemize} & 3 \\
\hline

US105 & En tant que \textbf{développeur}, je souhaite qu'un brouillon de solution soit généré à la fermeture d'un ticket afin de faciliter la documentation. &
\begin{itemize}[nosep, left=0pt, label={$\bullet$}]
\item Générer un brouillon de solution à partir du ticket via l'endpoint IA dédié, sans prétendre produire directement un article final publié.
\item Afficher ce brouillon dans l'interface du ticket pour permettre une relecture et une édition manuelle avant toute mise en forme documentaire.
\end{itemize} & 3 \\
\hline

US106 & En tant que \textbf{chef de projet}, je souhaite que l'IA propose une classification des demandes clients afin d'accélérer leur traitement. &
\begin{itemize}[nosep, left=0pt, label={$\bullet$}]
\item Utiliser l'endpoint IA de suggestions sur les demandes pour récupérer une classification automatique avec score de confiance.
\item Prévoir la reclassification et le feedback métier afin d'affiner les résultats de l'IA sur les demandes futures.
\end{itemize} & 3 \\
\hline

US107 & En tant que \textbf{développeur}, je souhaite que l'IA me recommande des articles de la base de connaissances afin de résoudre plus vite. &
\begin{itemize}[nosep, left=0pt, label={$\bullet$}]
\item Interroger l'endpoint IA de recommandations KB lié au ticket afin d'obtenir les articles de connaissance les plus pertinents.
\item Présenter ces suggestions directement dans la fiche ticket pour aider le développeur à résoudre plus vite.
\end{itemize} & 2 \\
\hline

US108 & En tant que \textbf{chef de projet}, je souhaite que l'IA structure automatiquement les solutions saisies en articles KB afin d'améliorer la documentation. &
\begin{itemize}[nosep, left=0pt, label={$\bullet$}]
\item S'appuyer sur le brouillon de solution généré par l'IA pour structurer une base de documentation réutilisable.
\item Laisser la possibilité de compléter, corriger et valider le contenu manuellement avant publication dans la base de connaissances.
\end{itemize} & 2 \\
\hline

US109 & En tant que \textbf{chef de projet}, je souhaite que l'IA fusionne le brouillon d'un blocage avec la solution humaine afin de produire un article KB définitif. &
\begin{itemize}[nosep, left=0pt, label={$\bullet$}]
\item Fusionner le brouillon de blocage et la solution issue du ticket pour produire un contenu de type base de connaissances prêt à relecture.
\item Conserver l'historique de l'analyse IA et de la résolution humaine afin d'éviter les doublons et de faciliter la validation finale.
\end{itemize} & 2 \\
\hline
\end{longtable}
\captionof{table}{Backlog du troisième sprint}\label{tab:backlog_sprint3_pro}
